\documentclass{article}
\usepackage[italian]{babel}
\usepackage[letterpaper,top=2cm,bottom=2cm,left=3cm,right=3cm,marginparwidth=1.75cm]{geometry}

\usepackage{mathtools}
\usepackage{amssymb}
\usepackage{amsmath}
\usepackage{amsthm}

\usepackage{graphicx}
\usepackage[colorlinks=true, allcolors=blue]{hyperref}
\title{Titolo}
\author{Autore}
\date{Data}
\begin{document}

\section*{Proprietà delle operazioni}

\subsection*{Proprietà associativa}

\begin{align*}
{\rm A} + ({\rm B} + {\rm C}) &= ({\rm A} + {\rm B}) + {\rm C}\\
{\rm A} \cdot ({\rm B} \cdot {\rm C}) &= ({\rm A} \cdot {\rm B}) \cdot {\rm C}.
\end{align*}

\paragraph{Mostra che sono uguali:}

\begin{align*}
a + (b + (c + d)) = (a + b) + (c + d) = ((a + b) + c) + d = (a + (b + c)) + d = a + ((b + c) + d).
\end{align*}

\paragraph{Mostra che sono uguali:}

\begin{align*}
a \cdot (b \cdot (c \cdot d)) = (a \cdot b) \cdot (c \cdot d) = ((a \cdot b) \cdot c) \cdot d = (a \cdot (b \cdot c)) \cdot d = a \cdot ((b \cdot c) \cdot d).
\end{align*}

\subsection*{Proprietà commutativa}

\begin{align*}
{\rm A} + {\rm B} &= {\rm B} + {\rm A}\\
{\rm A} \cdot {\rm B} &= {\rm B} \cdot {\rm A}.
\end{align*}

\paragraph{Mostra che sono uguali:}

\begin{align*}
a + b + c = a + c + b = c + a + b = c + b + a = b + c + a = b + a + c.
\end{align*}

\paragraph{Mostra che sono uguali:}

\begin{align*}
a \cdot b \cdot c = a \cdot c \cdot b = c \cdot a \cdot b = c \cdot b \cdot a = b \cdot c \cdot a = b \cdot a \cdot c.
\end{align*}

\subsection*{Proprietà distributiva}

\begin{align*}
{\rm A} \cdot ({\rm B} + {\rm C}) &= {\rm A} \cdot {\rm B} + {\rm A} \cdot {\rm C}.
\end{align*}

\paragraph{Mostra che:}

\begin{align*}
a \cdot (b + c + d) &= a \cdot b + a \cdot c + a \cdot d.
\end{align*}

\paragraph{Correggi l'errore:}

\begin{align*}
a = b &\longleftrightarrow a + c = b + c\\
a = b &\longleftrightarrow a \cdot c = b \cdot c\\
a < b &\longleftrightarrow a + c < b + c\\
a < b &\longleftrightarrow a \cdot c < b \cdot c.
\end{align*}

\newpage
\section*{Differenza e Divisione}

\begin{align*}
a - b = c &\longleftrightarrow a = b + c\\
a / b = c &\longleftrightarrow a = b \cdot c
\end{align*}

\paragraph{Mostra che:}

\begin{align*}
(a - c) + b = (a + b) - c\\ 
a + (b - c) = (a + b) - c\\
a - (c - b) = (a + b) - c\\
(a - b) - c = a - (b + c)\\
(a / c) \cdot b = (a \cdot b) / c\\ 
a \cdot (b / c) = (a \cdot b) / c\\
a / (c / b) = (a \cdot b) / c\\
(a / b) / c = a / (b \cdot c)\\
(a - b) \cdot c = a \cdot c - b \cdot c\\
(a + b) / c = (a / c) + (b / c)\\
(a - b) / c = (a / c) - (b / c).
\end{align*}

\paragraph{Correggi l'errore:}

\begin{align*}
a = b &\longleftrightarrow a - c = b - c\\
a = b &\longleftrightarrow c - a = c - b\\
a < b &\longleftrightarrow a - c < b - c\\
a < b &\longleftrightarrow c - a < c - b.
\end{align*}

\paragraph{Correggi l'errore:}

\begin{align*}
a = b &\longleftrightarrow a / c = b / c\\
a = b &\longleftrightarrow c / a = c / b\\
a < b &\longleftrightarrow a / c < b / c\\
a < b &\longleftrightarrow c / a < c / b.
\end{align*}

\newpage
\section*{Proprietà delle potenze}

\begin{align*}
{\rm A}^{{\rm B} + {\rm C}} &= {\rm A}^{\rm B} \cdot {\rm A}^{\rm C}\\
({\rm A}^{\rm B})^{\rm C} &= {\rm A}^{{\rm B} \cdot {\rm C}}\\
({\rm A} \cdot {\rm B})^{\rm C} &= {\rm A}^{\rm C} \cdot {\rm B}^{\rm C}.
\end{align*}

\paragraph{Mostra che:}

\begin{align*}
a^{b + c + d} &= a^b \cdot a^c \cdot a^d\\
((a^b)^c)^d &= a^{b \cdot c \cdot d}\\
(a \cdot b \cdot c)^d &= a^d \cdot b^d \cdot c^d.
\end{align*}

\paragraph{Correggi l'errore:}

\begin{align*}
a = b &\longleftrightarrow a^c = b^c\\
a = b &\longleftrightarrow c^a = c^b\\
a < b &\longleftrightarrow a^c < b^c\\
a < b &\longleftrightarrow c^a < c^b.
\end{align*}

\newpage
\section*{Radicali e Logaritmi}

\begin{align*}
\sqrt[b]{a} = c &\longleftrightarrow a = c^b\\
\log_b{a} = c &\longleftrightarrow a = b^c.
\end{align*}

\paragraph{Mostra che:}

\begin{align*}
a^{b - c} = a^b / a^c\\
a^{b / c} = \sqrt[c]{(a^b)}\\
a^{b / c} = (\sqrt[c]{a})^b\\
\sqrt[c]{a \cdot b} = \sqrt[c]{a} \cdot \sqrt[c]{b}\\
(a / b)^c = a^c / b^c\\
\sqrt[c]{(a / b)} = \sqrt[c]{a} / \sqrt[c]{b.}
\end{align*}

\paragraph{Mostra che:}

\begin{align*}
\log_a(b \cdot c) = \log_a(b) + \log_a(c)\\
\log_a(b / c) = \log_a(b) - \log_a(c)\\
\log_a(b^c) = c \cdot \log_a(b)\\
\log_a(\sqrt[c]{b}) = \log_a(b) / c\\
\log_a(b) \cdot \log_b(c) = \log_a(c)\\
\log_a(b) = \log_c(b) / \log_c(a)
\end{align*}

\paragraph{Correggi l'errore:}

\begin{align*}
a = b &\longleftrightarrow \sqrt[c]{a} = \sqrt[c]{b}\\
a = b &\longleftrightarrow \sqrt[a]{c} = \sqrt[b]{c}\\
a < b &\longleftrightarrow \sqrt[c]{a} < \sqrt[c]{b}\\
a < b &\longleftrightarrow \sqrt[a]{c} < \sqrt[b]{c}
\end{align*}

\paragraph{Correggi l'errore:}

\begin{align*}
a = b &\longleftrightarrow \log_c(a) = \log_c(b)\\
a = b &\longleftrightarrow \log_a(c) = \log_b(c)\\
a < b &\longleftrightarrow \log_c(a) < \log_c(b)\\
a < b &\longleftrightarrow \log_a(c) < \log_b(c)
\end{align*}

\newpage
\section*{Aritmetica}

\subsection*{Massimi e Minimi}

\begin{align*}
c \subseteq (a \cap b) &\longleftrightarrow (c \subseteq a) \land (c \subseteq b)\\
(a \cup b) \subseteq c &\longleftrightarrow (a \subseteq c) \land (b \subseteq c)\\
c \le \min(a, b) &\longleftrightarrow (c \le a) \land (c \le b)\\
\max(a, b) \le c &\longleftrightarrow (a \le c) \land (b \le c)\\
c \le \max(a, b) &\longleftrightarrow (c \le a) \lor (c \le b)\\
\min(a, b) \le c&\longleftrightarrow (a \le c) \lor (b \le c)
\end{align*}

\paragraph{Mostra che:}

\begin{align*}
a \subseteq b &\longrightarrow \min(a) > \min(b)\\
a \subseteq b &\longrightarrow \max(a) < \max(b)
\end{align*}

\paragraph{Mostra che:}

\begin{align*}
\max(a, b) + \min(a, b) = a + b
\end{align*}

\subsection*{Divisione Euclidea}

\begin{align*}
a = (a // b) \cdot b + a \% b\\
a \% b < b\\
(x < b \land a = x \cdot b + y) \rightarrow (x = a // b) \land (y = a \% b) 
\end{align*}

\paragraph{Mostra che:}

\begin{align*}
((b \% a) + (c \% a)) \% a &= (b + c) \% a\\
((b \% a) \cdot (c \% a)) \% a &= (b \cdot c) \% a
\end{align*}

\paragraph{Mostra che:} se $a \% b = 0$

\begin{align*}
(c \% a) \% b &= c \% b\\
c // b &= (c // a) \cdot (a // b) + ((c \% a) // b)\\
(y < a // b \land z < b \land c = x \cdot a + y \cdot b + z) &\rightarrow x, y, z = c // a, (c \% a) // b, c \% b
\end{align*}

\subsection*{Fattoriale}

\begin{align*}
(n + 1)! = (n + 1) \cdot n!
\end{align*}

\paragraph{Mostra che:}

\begin{align*}
\frac{(a + b - 1)!}{(a - 1)! \cdot b!} + \frac{(a + b - 1)!}{a! \cdot (b - 1)!} = \frac{(a + b)!}{a! \cdot b!}
\end{align*}

\newpage
\section*{Fattorizzazione unica}

\paragraph{Numeri primi e irriducibili}

\begin{align*}
p = a \cdot b &\longrightarrow (a^2 - 1) \cdot (b^2 - 1) = 0\\
(a \cdot b) \% p = 0 &\longrightarrow (a \% p) \cdot (b \% p) = 0
\end{align*}

\paragraph{Definizione di gcd ed lcm}

\begin{align*}
\operatorname{gcd}(a, b) \% c = 0 &\longleftrightarrow (a \% c = 0) \land (b \% c = 0)\\
c \% \operatorname{lcm}(a, b) = 0 &\longleftrightarrow (c \% a = 0) \land (c \% b = 0)
\end{align*}

\paragraph{Se per ogni primo p:}

\begin{align*}
\max_{a \% p^d = 0}(d) = \max_{b \% p^d = 0}(d),
\end{align*}

allora $a = b$.

\paragraph{Mostra che:} sia $p$ primo

\begin{align*}
\max_{(a \cdot b) \% p^d = 0}(d) &= \max_{a \% p^d = 0}(d) + \max_{b \% p^d = 0}(d)
\end{align*}

\paragraph{Mostra che:} sia $p$ primo

\begin{align*}
\max_{\operatorname{lcm}(a, b) \% p^d = 0}(d) &= \max\left(\max_{a \% p^d = 0}(d), \max_{b \% p^d = 0}(d)\right)\\
\max_{\operatorname{gcd}(a, b) \% p^d = 0}(d) &= \min\left(\max_{a \% p^d = 0}(d), \max_{b \% p^d = 0}(d)\right)
\end{align*}

\paragraph{Mostra che:}

\begin{align*}
\operatorname{lcm}(a, b) \cdot \operatorname{gcd}(a, b) = a \cdot b
\end{align*}

\newpage
\section*{Insiemi e logica}

\subsection*{Tavole di verità}

\begin{align*}
\lnot 0 = 1\\
\lnot 1 = 0\\
1 \land 1 = 1 \lor 1 = 1\\
0 \land 0 = 0 \lor 0 = 0\\
0 \land 1 = 1 \land 0 = 0\\
0 \lor 1 = 1 \lor 0 = 1.
\end{align*}

\paragraph{Dimostra ciascuna delle seguenti utilizzando le altre:}

\begin{align*}
a &= \lnot (\lnot a)\\
a &= a \land a\\
a &= a \lor a\\
\lnot (a \land b) &= (\lnot a) \lor (\lnot b)\\
\lnot (a \lor b) &= (\lnot a) \land (\lnot b)\\
a \land (b \lor c) &= (a \land b) \lor (a \land c)\\
a \lor (b \land c) &= (a \lor b) \land (a \lor c).
\end{align*}

\subsection*{Insiemi}

\begin{align*}
{\rm A} \setminus {\rm B} &= \{a: a \in {\rm A} \land a \notin {\rm B}\}\\
{\rm A} \cap {\rm B} &= \{a: a \in {\rm A} \land a \in {\rm B}\}\\
{\rm A} \cup {\rm B} &= \{a: a \in {\rm A} \lor a \in {\rm B}\}.
\end{align*}

\paragraph{Mostra che valgono le seguenti:}

\begin{align*}
(a \setminus b) \setminus c &= (a \setminus b) \cap (a \setminus c)\\
a \setminus (b \cup c) &= (a \setminus b) \cap (a \setminus c)\\
a \cap (b \setminus c) &= (a \cap b) \setminus (a \cap c)\\
(a \cap b) \setminus c &= (a \cap b) \setminus (a \cap c)\\
a \setminus (b \setminus c) &= (a \cap c) \cup (a \setminus b)\\
a \setminus (b \cap c) &= (a \setminus b) \cup (a \setminus b)\\
(a \cup b) \setminus c &= (a \setminus c) \cup (a \setminus c)\\
a \cup (b \setminus c) &= (a \cup b) \setminus (c \setminus a).
\end{align*}

\newpage
\section*{Funzioni}

\paragraph{Una funzione $f: {\rm A} \rightarrow {\rm B}$ si dice:}

\begin{itemize}
\item Sugettiva, se $\{f(a): a \in {\rm A}\} = {\rm B}$.
\item Iniettiva, se $f(a) = f(b) \longrightarrow a = b$.
\item Di periodo $a$, se $f(x + a) = f(x)$.
\item Pari, se $f(- x) = f(x)$.
\item Dispari, se $f(- x) = - f(x)$.
\end{itemize}

\paragraph{Non si fa mai:}

\begin{itemize}
\item Dividere per $0$, fare il logaritmo in base $1$, fare $0^0$.
\item Fare la radice pari di un numero negativo.
\item Elevare ad un esponente denominatore pari un numero negativo.
\item Elevare ad un esponente irrazionale un numero negativo.
\item L'esponenziale con base negativa.
\item Il logaritmo con base negativa.
\item Il logaritmo con argomento negativo.
\end{itemize}

\paragraph{Mostra che:}

\begin{itemize}
\item Composizione di funzioni surgettive è surgettiva.
\item Composizione di funzioni iniettive è iniettiva.
\item Se $f \circ g$ è surgettiva, $f$ è sugettiva.
\item Se $f \circ g$ è iniettiva, $g$ è iniettiva.
\item Se $f \circ g$ è surgettiva ed $f$ è iniettiva, $g$ è surgettiva.
\item Se $f \circ g$ è iniettiva e $g$ è surgettiva, $f$ è iniettiva.
\item Composizione di funzioni pari è pari.
\item Composizione di funzioni dispari è pari.
\item Composizione di una funzione pari e di una dispari è dispari.
\item Composizione di una funzione dispari e di una pari è dispari.
\item Se $g$ è di periodo $a$, $f \circ g$ è di periodo $\rm T$.
\item Se $f$ è di periodo $a$ e $b$, allora è di periodo $c \cdot \gcd(a, b)$.
\item Una potenza con base positiva è positiva e monotona nella base e nell'esponente.
\item Un logaritmo con base ed argomento positivi è monotono nella base e nell'argomento.
\end{itemize}

\paragraph{Determina dominio ed immagine delle seguenti funzioni:}

\begin{align*}
x + a \qquad a + x \qquad x - a \qquad a - x \qquad x \cdot a \qquad a \cdot x \qquad x / a\\
a / x \qquad x^a \qquad a^x \qquad \sqrt[a]{x} \qquad \sqrt[x]{a} \qquad \log_a(x) \qquad \log_x(a).
\end{align*}

\paragraph{Determina poi se sono pari, dispari, periodiche, monotone, iniettive, surgettive.}

\newpage
\section*{Condizioni di esistenza}

\paragraph{Data una espressione} $f(x)$:

\begin{itemize}
\item è possibile determinare il dominio naturale di $f$ imponendo le condizioni di esistenza.
\item è possibile risolvere $f(x) = y$ per alcuni $y$, tali valori sono detti immagine di $f$.
\item si ottiene poi $x = g(y)$ se $f$ è iniettiva, in tal caso $g$ ha come dominio l'immagine di $f$ e viceversa.
\item altrimenti si può ottenere $x \in \{g_0(y), \dots, g_k(y)\}$. In tal caso:
\begin{itemize}
\item si restringe il dominio di $g_h$ all'insieme delle $y$ per cui $f (g_h (y)) = y$
\item l'unione dei domini di $g_0, \dots, g_k$ è l'immagine di $f$
\item l'immagine di $g_h$ è $x = g_h (f (x))$ per $x$ nel dominio di $f$
\item le immagini di $g_0, \dots, g_k$ sono una partizione del dominio di $f$
\item restringendo il dominio di $f$ all'immagine di $g_h$, $f$ risulta bigettiva sul dominio di $g_h$.
\end{itemize}
\end{itemize}

\paragraph{Esempio:} $f(x) = \sqrt{(x - a_0)(x - a_1)}$ ha come dominio $(\infty, \min(a_0, a_1)] \cup [\max(a_0, a_1), \infty)$.

\begin{align*}
x^2 - (a_0 + a_1) x + a_0 a_1 - y^2 = 0
\end{align*}

ha come soluzioni

\begin{align*}
g_0(y) = \frac{a_0 + a_1}{2} - \sqrt{\frac{(a_0 - a_1)^2}{4} + y^2}\\
g_1(y) = \frac{a_0 + a_1}{2} + \sqrt{\frac{(a_0 - a_1)^2}{4} + y^2.}
\end{align*}

\begin{itemize}
\item Entrambe ammettono come dominio $\mathbb R$, ma ristretto a $\sqrt{y^2} = y$, cioé $[0, \infty)$.
\item L'immagine di $g_0$ è $(\infty, \min(a_0, a_1)]$, mentre quella di $g_1$ è $[\max(a_0, a_1), \infty)$.
\end{itemize}

\paragraph{Esempio:} $f(x) = \sqrt{- (x - a_0)(x - a_1)}$ ha come dominio $[\min(a_0, a_1), \max(a_0, a_1)]$.

\begin{align*}
x^2 - (a_0 + a_1) x + a_0 a_1 + y^2 = 0
\end{align*}

ha come soluzioni

\begin{align*}
g_0(y) = \frac{a_0 + a_1}{2} - \sqrt{\frac{(a_0 - a_1)^2}{4} - y^2}\\
g_1(y) = \frac{a_0 + a_1}{2} + \sqrt{\frac{(a_0 - a_1)^2}{4} - y^2.}
\end{align*}

\begin{itemize}
\item Entrambe ammettono come dominio $\left[0, \frac{|a_0 - a_1|}{2}\right]$.
\item L'immagine di $g_0$ è $\left[\min(a_0, a_1), \frac{a_0 + a_1}{2}\right]$, mentre quella di $g_1$ è $\left[\frac{a_0 + a_1}{2}, \max(a_0, a_1)\right]$.
\end{itemize}

\paragraph{Esempio:} $f(x) = \sqrt{(x - a)^2 + b^2}$ ha come dominio $\mathbb R$.

\begin{align*}
g_0(y) = a - \sqrt{y^2 - b^2}\\
g_1(y) = a + \sqrt{y^2 - b^2.}
\end{align*}

\begin{itemize}
\item Entrambe ammettono come dominio $(\infty, |b|] \cup [|b|, \infty)$, ma ristretto a $\sqrt{y^2} = y$, cioé $[|b|, \infty)$.
\item L'immagine di $g_0$ è $(\infty, a]$, mentre quella di $g_1$ è $[a, \infty)$.
\end{itemize}

\newpage
\section*{Geometria}

\paragraph{Trasformazioni conformi:}

\begin{itemize}
\item Dati due punti ${\rm A}, {\rm B}$, esiste un solo vettore $u$ tale che ${\rm A} + u = {\rm B}$.
\item Dati due vettori $u, v$ paralleli, esiste un solo $x$ per cui $v = x \cdot u$.
\item Dati due vettori $u, v$ di stesso modulo, esiste un solo $\alpha$ per cui $v = r_\theta \cdot u$.
\item Dati due vettori $u, v$, esistono unici $a, b$ tali che $v = a \cdot u + b \cdot (r_\frac{\pi}{2} \cdot u)$.
\end{itemize}

\paragraph{Proprietà:}

\begin{itemize}
\item $({\rm A} + u) + v = {\rm A} + (u + v) $
\item $r_\alpha \cdot (r_\beta \cdot u) = r_{\alpha + \beta} \cdot u$
\item $x \cdot (y \cdot u) = (y \cdot x) \cdot u$
\item $r_\alpha \cdot (x \cdot u) = x \cdot (r_\alpha \cdot u)$
\item $r_\alpha \cdot (u + v) = r_\alpha \cdot u + r_\alpha \cdot v$
\item $x \cdot (u + v) = x \cdot u + x \cdot v$
\item $(x + y) \cdot u = x \cdot u + y \cdot u$
\item $|x \cdot u| = |x| \cdot |u|$
\item $r_\pi \cdot u = (- 1) \cdot u$
\end{itemize}

\paragraph{Mostra che:}

\begin{itemize}
\item Una traslazione composta una rotazione è una rotazione composta una traslazione.
\item Una traslazione composta una dilatazione è una dilatazione composta una traslazione.
\item Le traslazioni mandano rette in rette parallele.
\item Le traslazioni conservano le distanze.
\item Le rotazioni conservano gli angoli.
\item Le rotazioni conservano le distanze.
\item Le dilatazioni mandano rette in rette parallele.
\item Le dilatazioni moltiplicano le distanze per una costante.
\item Dati due segmenti orientati esiste una unica trasformazione conforme che manda l'uno nell'altro.
\end{itemize}

\newpage
\section*{Prodotto scalare e vettoriale}

\paragraph{Siano $v = a \cdot u + b \cdot r_\frac{\pi}{2} \cdot u$:}

\begin{itemize}
\item Si dice prodotto scalare $u \cdot v = a \cdot |u|^2$.
\item Si dice prodotto vettoriale $u \times v = b \cdot |u|^2$.
\end{itemize}

\paragraph{Mostra che:}

\begin{align*}
u \cdot u = |u|^2 &\qquad u \times u = 0\\
u \cdot v = v \cdot u &\qquad u \times v + v \times u = 0\\
u \cdot (r_\frac{\pi}{2} \cdot u) = 0 &\qquad u \times (r_\frac{\pi}{2} \cdot u) = |u|^2\\
(r_\alpha \cdot u) \cdot (r_\alpha \cdot v) = u \cdot v &\qquad (r_\alpha \cdot u) \times (r_\alpha \cdot v) = u \times v\\
u \cdot (x \cdot v + w) = x \cdot (u \cdot v) + (u \cdot w) &\qquad u \times (x \cdot v + w) = x \cdot (u \times v) + (u \times w).
\end{align*}

\paragraph{Mostra che:}

\begin{align*}
u \times v &= (r_\frac{\pi}{2} \cdot u) \cdot v\\
|u \cdot v|^2 + |u \times v|^2 &= |u|^2 \cdot |v|^2.    
\end{align*}

\subsection*{Seno e Coseno}

\paragraph{Siano $r_\alpha \cdot u = a \cdot u + b \cdot (r_\frac{\pi}{2} \cdot u)$:}

\begin{itemize}
\item Si dice coseno di $\alpha$ il valore $a = \cos(\alpha)$.
\item Si dice seno di $\alpha$ il valore $b = \sin(\alpha)$.
\end{itemize}

\paragraph{Mostra che:}

\begin{align*}
\cos(\alpha + \beta) &= \cos(\alpha) \cdot \cos(\beta) - \sin(\alpha) \cdot \sin(\beta)\\
\sin(\alpha + \beta) &= \sin(\alpha) \cdot \cos(\beta) + \cos(\alpha) \cdot \sin(\beta).
\end{align*}

\paragraph{Mostra che:}

\begin{align*}
\cos(\alpha) \cdot \cos(\beta) &= \frac{\cos(\alpha + \beta) + \cos(\alpha - \beta)}{2}\\
\sin(\alpha) \cdot \cos(\beta) &= \frac{\sin(\alpha + \beta) + \sin(\alpha - \beta)}{2}\\
\sin(\alpha) \cdot \sin(\beta) &= \frac{\cos(\alpha - \beta) - \cos(\alpha - \beta)}{2.}
\end{align*}

\paragraph{Mostra che:}

\begin{align*}
\cos(\alpha) + \cos(\beta) &= 2 \cdot \left(\cos\left(\frac{\alpha + \beta}{2}\right) \cdot \cos\left(\frac{\alpha - \beta}{2}\right)\right)\\
\cos(\alpha) - \cos(\beta) &= 2 \cdot \left(\sin\left(\frac{\alpha + \beta}{2}\right) \cdot \sin\left(\frac{\alpha - \beta}{2}\right)\right)\\
\sin(\alpha) + \sin(\beta) &= 2 \cdot \left(\sin\left(\frac{\alpha + \beta}{2}\right) \cdot \cos\left(\frac{\alpha - \beta}{2}\right)\right).
\end{align*}

\newpage
\section*{Coniche}

\paragraph{L'equazione di una parabola} di fuoco $(\alpha, \beta)$ e direttrice $\{y = m x + q\}$ è:

\begin{align*}
\frac{(x + m y)^2}{m^2 + 1} - 2 \left(\alpha + \frac{m q}{m^2 + 1}\right) x - 2 \left(\beta - \frac{q}{m^2 + 1}\right) y + \left(\alpha^2 + \beta^2 - \frac{q^2}{m^2 + 1}\right) = 0.
\end{align*}

Per ogni punto della parabola passa una sola retta tangente.

\begin{itemize}
\item Per i punti esterni vale $>$ e passano $2$ rette tangenti.
\item Per i punti interni vale $<$ e non passano rette tangenti.
\end{itemize}

\paragraph{L'equazione di una conica} di fuoco $(\alpha, \beta)$, direttrice $\{y = m x + q\}$, ed eccentricità $\varepsilon$ è:

\begin{align*}
x^2 + y^2 - \frac{\varepsilon (m x - y)^2}{m^2 + 1} - 2 \left(\alpha + \frac{\varepsilon m q}{m^2 + 1}\right) x - 2 \left(\beta - \frac{\varepsilon q}{m^2 + 1}\right) y + \left(\alpha^2 + \beta^2 - \frac{\varepsilon q^2}{m^2 + 1}\right) = 0.
\end{align*}

Per ogni punto della conica passa una sola retta tangente.

\begin{itemize}
\item Per i punti esterni vale $>$ e passano $2$ rette tangenti.
\item Per i punti interni vale $<$ e non passano rette tangenti.
\end{itemize}

\paragraph{Dati due punti:} $(\alpha_0, \beta_0), (\alpha_1, \beta_1)$ diciamo

\begin{align*}
4 c^2 = (\alpha_0 - \alpha_1)^2 + (\beta_0 - \beta_1)^2
\end{align*}

L'equazione di un'ellisse di asse $a > c$ o di un'iperbole di asse $a < c$ è

\begin{align*}
4 a^2 \left((x - \alpha)^2 + (y - \beta)^2 + c^2 - a^2\right) = ((x - \alpha) (\alpha_1 - \alpha_0) + (y - \beta) (\beta_1 - \beta_0))^2.
\end{align*}

\begin{itemize}
\item Per i punti esterni vale $>$ e passano $2$ rette tangenti.
\item Per i punti interni vale $<$ e non passano rette tangenti.
\end{itemize}

Nel caso $a = c$ l'equazione si riduce a quella di una retta passante per i due fuochi.

\begin{itemize}
\item Se $\beta_1, \beta_0, \alpha = 0$ abbiamo $\frac{x^2}{a^2} + \frac{y^2}{a^2 - c^2} = 1$
\item Se $\alpha_1, \alpha_0, \beta = 0$ abbiamo $\frac{x^2}{a^2 - c^2} + \frac{y^2}{a^2} = 1$
\end{itemize}

\newpage
\section*{Circonferenze}

\paragraph{L'equazione di una circonferenza} di centro $(\alpha, \beta)$ e raggio $\rho$ è:

\begin{align*}
x^2 + y^2 - 2 \alpha x - 2 \beta y + (\alpha^2 + \beta^2 - \rho^2) = 0.
\end{align*}

Per ogni punto della circonferenza passa una sola retta tangente.

\begin{itemize}
\item Per i punti esterni vale $x^2 + y^2 - 2 \alpha x - 2 \beta y + (\alpha^2 + \beta^2 - \rho^2) > 0$ e passano $2$ rette tangenti.
\item Per i punti interni vale $x^2 + y^2 - 2 \alpha x - 2 \beta y + (\alpha^2 + \beta^2 - \rho^2) < 0$ e non passano rette tangenti.
\end{itemize}

\paragraph{Due circonferenze} di raggi $\rho_1, \rho_0$ e centri distanti $\delta$ possono essere:

\begin{itemize}
\item Esterne se $\delta > \rho_1 + \rho_0$ (intersezione vuota).
\item Tangenti esternamente se $\delta = \rho_1 + \rho_0$ (un punto di intersezione).
\item Secanti se $|\rho_1 - \rho_0| < \delta < \rho_1 + \rho_0$ (due punti di intersezione).
\item Tangenti internamente se $\delta = |\rho_1 - \rho_0|$ (un punto di intersezione).
\item Interne se $\delta < |\rho_1 - \rho_0|$ (intersezione vuota).
\end{itemize}

I punti di intersezione si trovano intersecando una delle due circonferenze con la retta di equazione

\begin{align*}
2 (\alpha_1 - \alpha_0) x + 2 (\beta_1 - \beta_0) y = (\alpha_1^2 + \beta_1^2 - \rho_1^2) - (\alpha_0^2 + \beta_0^2 - \rho_0^2)
\end{align*}

\paragraph{Data una retta generica} $r = \{y = m x + q\}$:

\begin{itemize}
\item $r$ è secante ad una circonferenza se $(m \alpha + q - \beta)^2 < \rho^2 (m^2 + 1)$
\item $r$ è tangente ad una circonferenza se $(m \alpha + q - \beta)^2 = \rho^2 (m^2 + 1)$
\item $r$ è disgiunta da una circonferenza se $(m \alpha + q - \beta)^2 > \rho^2 (m^2 + 1)$.
\end{itemize}

\subsection*{Esercizio}

\paragraph{Data l'equazione di una circonferenza $\cal C$:}

\begin{itemize}
\item Determina il centro $(\alpha, \beta)$ ed il raggio $\rho$ di $\cal C$.
\item Dato un punto ${\rm P}$, determina se $\rm P$ appartiene a $\cal C$, se è interno o esterno. Se $\rm P$ è esterno:
\begin{itemize}
\item determina le tangenti a $\cal C$ passanti per $\rm P$
\item queste intersecano $\cal C$ in in punto ciascuna, determina le loro coordinate
\item determina l'equazione della retta passante per i due punti trovati.
\end{itemize}
\item Data una retta $r$, determina se $r$ è disgiunta, secante o tangente a $\cal C$. Se $r$ è secante:
\begin{itemize}
\item determina i punti di intersezione di $r$ con $\cal C$
\item per ciascuno di questi punti passa una retta tangente a $\cal C$, determina le loro equazioni
\item determina le coordinate del punto di intersezione delle rette trovate.
\end{itemize}
\item Dati $a, b$, per quali $c$ la circonferenza $\{x^2 + y^2 + a x + b y + c = 0\}$ e la circonferenza $\cal C$ sono interne (e quale delle due contiene l'altra), esterne o secanti (e determinare i punti di intersezione).
\end{itemize}

\newpage
\section*{Limiti e continuità}

\paragraph{Si dice che $\lim_{n \to \infty} a_n = x$ se:}

\begin{align*}
\forall \varepsilon > 0, \quad \exists n \in \mathbb N, \quad \forall m > n, \quad |a_n - x| < \varepsilon.
\end{align*}

\paragraph{Si dice che $\lim_{y \rightarrow x} f(y) = z$ se:}

\begin{align*}
\forall \lim_{n \rightarrow \infty} a_n = x, \quad \lim_{n \rightarrow \infty} f(a_n) = z.
\end{align*}

\paragraph{La chiusura di un insieme consiste dei limiti di successioni a valori in quell'insieme.}

\paragraph{Un insieme si dice aperto se ogni successione che non lo interseca non ci converge.}

\paragraph{Si dice che una funzione $f$ è continua in $x$ se:}

\begin{align*}
\lim_{y \rightarrow x} f(y) = f(x).
\end{align*}

\paragraph{Una funzione continua manda intervalli in intervalli.}

\paragraph{Una funzione continua manda chiusi limitati in chiusi limitati.}

\subsection*{Esercizi}

\paragraph{Sia $a_n$ una successione crescente, mostra che:}

\begin{align*}
\lim_{n \to \infty} a_n = \sup_{n \in {\mathbb N}} a_n.
\end{align*}

\paragraph{Sia $a_n$ una successione decrescente, mostra che:}

\begin{align*}
\lim_{n \to \infty} a_n = \inf_{n \in {\mathbb N}} a_n.
\end{align*}

\paragraph{Sia $a_n$ una successione, mostra che sono monotone:}

\begin{align*}
\left(\sup_{m > n} a_m\right)_{n \in {\mathbb N}} \qquad \left(\inf_{m > n} a_m\right)_{n \in {\mathbb N}.}
\end{align*}

\paragraph{Sia $a_n$ una successione convergente, mostra che:}

\begin{align*}
\lim_{n \to \infty} a_n &= \inf_{n \in {\mathbb N}}\left(\sup_{m > n} a_m\right)\\
\lim_{n \to \infty} a_n &= \sup_{n \in {\mathbb N}}\left(\inf_{m > n} a_m\right).
\end{align*}

\paragraph{Mostra che $\lim_{y \rightarrow x} f(y) = z$ se e solo se:}

\begin{align*}
\forall \varepsilon > 0, \quad \exists \delta > 0, \quad \forall |y - x| < \delta, \quad |f(y) - z| < \varepsilon.
\end{align*}

\paragraph{Mostra che $f$ è una funzione continua in $x$ se e solo se:}

\begin{align*}
\lim_{n \to \infty} a_n = x \longrightarrow \lim_{n \to \infty} f(a_n) = f(x).
\end{align*}

\paragraph{Mostra che:}

\begin{itemize}
\item Una funzione è continua se e solo se chiusura di immagine contiene immagine di chusura.
\item Una funzione è continua se e solo se controimmagine di un aperto è aperta.
\end{itemize}

\newpage
\section*{Derivate}

\paragraph{Si dice che una funzione $f$ è derivabile in $x$ se esiste il limite del rapporto incrementale:}

\begin{align*}
\frac{\partial f(y)}{\partial y}(x) = \lim_{y \rightarrow x} \frac{f(y) - f(x)}{y - x.}
\end{align*}

\paragraph{Si dice che una funzione $f$ è crescente se:}

\begin{align*}
\forall x < y, \quad f(x) < f(y).
\end{align*}

\paragraph{Si dice che una funzione $f$ è convessa se:}

\begin{align*}
\forall x < y < z, \quad f(y) < \left(\frac{z - y}{z - x}\right) \cdot f(x) + \left(\frac{y - x}{z - x}\right) \cdot f(z).
\end{align*}

\subsection*{Esercizi}

\paragraph{Mostra che una funzione $f$ crescente se e solo se:}

\begin{align*}
\forall x < y, \quad \frac{f(y) - f(x)}{y - x} > 0.
\end{align*}

\paragraph{Mostra che una funzione è crescente se e solo se il suo rapporto incrementale è positivo.}

\paragraph{Mostra che una funzione è crescente se e solo se la sua derivata è positiva.}

\paragraph{Mostra che una funzione $f$ è convessa se e solo se:}

\begin{align*}
\forall x < y < z, \quad \frac{f(y) - f(x)}{y - x} < \frac{f(z) - f(y)}{z - y.}
\end{align*}

\paragraph{Mostra che una funzione è convessa se e solo se la sua derivata è crescente.}

\paragraph{Mostra che una funzione è convessa se e solo se il suo rapporto incrementale è monotono.}

\paragraph{Mostra che:}

\begin{align*}
\partial (f \cdot g) &= (\partial f) \cdot g + f \cdot (\partial g)\\
\partial (f \circ g) &= ((\partial f) \circ g) \cdot \partial g\\
\partial (f^{- 1}) &= 1 / ((\partial f) \circ (f^{- 1})).
\end{align*}

\newpage
\section*{Limiti notevoli}

\paragraph{Si riesce a calcolare che:}

\begin{align*}
\lim_{x \rightarrow 0}\frac{\sin(x)}{x} &= 1\\
\lim_{x \rightarrow 0}\frac{\cos(x) - 1}{x^2} &= \frac{1}{2}\\
\lim_{x \rightarrow 0}\frac{\exp(x) - 1}{x} &= 1\\
\lim_{x \rightarrow 0}\frac{\log(x + 1)}{x} &= 1
\end{align*}

\paragraph{Si riesce a calcolare che:}

\begin{align*}
\frac{\partial \sin(y)}{\partial y}(x) &= \cos(x)\\
\frac{\partial \cos(y)}{\partial y}(x) &= - \sin(x)\\
\frac{\partial \exp(y)}{\partial y}(x) &= \exp(x)\\
\frac{\partial \log(y)}{\partial y}(x) &= \frac{1}{x}\\
\frac{\partial y^a}{\partial y}(x) &= a \cdot x^{a - 1}.
\end{align*}

\paragraph{Si definiscono:}

\begin{align*}
\tan(x) &= \frac{\sin(x)}{\cos(x)}\\
\cot(x) &= \frac{\cos(x)}{\sin(x).}
\end{align*}

\subsection*{Esercizi}

\paragraph{Mostra che sono bigettive:}

\begin{align*}
\cos:& [0, \pi] \rightarrow [- 1, 1]\\
\sin:& \left[- \frac{\pi}{2}, \frac{\pi}{2}\right] \rightarrow [- 1, 1]\\
\tan:& \left(- \frac{\pi}{2}, \frac{\pi}{2}\right) \rightarrow {\mathbb R}\\
\cot:& (0, \pi) \rightarrow {\mathbb R}
\end{align*}

\paragraph{Determina poi se sono pari, dispari, periodiche, monotone.}

\paragraph{Determina poi se le loro inverse sono pari, dispari, periodiche, monotone.}

\newpage
\section*{Funzioni iperboliche}

\paragraph{Si definiscono:}

\begin{align*}
\cosh(x) &= \frac{\exp(x) + \exp(- x)}{2}\\
\sinh(x) &= \frac{\exp(x) - \exp(- x)}{2.}
\end{align*}

\paragraph{Si definiscono:}

\begin{align*}
\tanh(x) &= \frac{\sinh(x)}{\cosh(x)}\\
\coth(x) &= \frac{\cosh(x)}{\sinh(x).}
\end{align*}

\subsection*{Esercizi}

\paragraph{Mostra che sono bigettive:}

\begin{align*}
\cosh:& {\mathbb R} \rightarrow {\mathbb R}_+\\
\sinh:& {\mathbb R} \rightarrow {\mathbb R}\\
\tanh:& {\mathbb R} \rightarrow [- 1, 1]\\
\coth:& {\mathbb R} \setminus \{0\} \rightarrow {\mathbb R} \setminus [-1, 1].
\end{align*}

\paragraph{Determina poi se sono pari, dispari, periodiche, monotone.}

\paragraph{Calcola le loro inverse e determina se sono pari, dispari, periodiche, monotone.}

\end{document}